\documentclass[11pt]{amsart}
\usepackage[margin=1.15in]{geometry}
\usepackage{amsmath,amssymb}
\usepackage[hidelinks]{hyperref}

\newcommand{\CF}{\mathrm{CF}}

\title{Compact Curve-Flat Quotient Realization: Literature Status}
\author{Agent lane 10}
\date{2026-06-18}

\begin{document}
\maketitle

\section*{Status}

Status: \texttt{literature\_already\_answered}.

\section*{Source Question}

In Remark 6.9 of Flores--Jung--Lancien--Petitjean--Prochazka--Quilis,
\emph{On curve-flat Lipschitz functions and their linearizations}
\cite{FJLPPQ2024}, the authors give a compact example whose
curve-flat quotient is the interval $[0,1]$. They then propose the
following compact realization problem: study for which non-purely
$1$-unrectifiable compact metric spaces $K$ there exists a compact
metric space $M$ such that
\[
        M_{\CF}=K .
\]

\section*{Supporting Answer}

Kaasik--Quilis \cite{KQ2026} explicitly cite this question in their
introduction. They explain that their exact complete-space theorem
does not settle the compact version because the constructed space is
not compact unless the target is finite, and then state that they
modify the construction in order to solve the compact problem, with
the caveat that exact isometry is not always obtained.

The answer supplied by \cite{KQ2026} is therefore scoped as follows.
Their Theorem B states that for every compact metric space $(M,d)$,
every countable ordinal $\alpha$, and every $\varepsilon>0$, there is a
compact metric space $(N,\rho)$ with $\alpha_{cf}(N)\geq\alpha$ such
that
\[
        N/\rho_{cf}^{\alpha}
\]
is $(1+\varepsilon)$-isometric to $(M,d)$. The more technical Theorem
3.1 gives exact isometry in particular when the compact target is
finite, geodesic, or a gapped segment.

\section*{Scope Limitations}

This is not recorded as a new proof by the run. It is a separate
later-paper answer to the source paper's compact curve-flat quotient
realization program. The all-compact statement is approximate up to
arbitrarily small bi-Lipschitz distortion. The packet does not assert
that every compact metric space has an exact compact first-order
curve-flat preimage.

\section*{Citation Evidence}

The supporting paper explicitly names the question from
\cite{FJLPPQ2024}: it says that Theorem A does not answer that question
because the constructed space need not be compact, and then introduces
the compact modification leading to Theorem B. Thus this belongs in the
literature-already-answered folder, not in a new full or partial proof
folder.

\begin{thebibliography}{9}

\bibitem{FJLPPQ2024}
G. Flores, M. Jung, G. Lancien, C. Petitjean, A. Prochazka, and
A. Quilis,
\emph{On curve-flat Lipschitz functions and their linearizations},
arXiv:2411.08369.

\bibitem{KQ2026}
J. K. Kaasik and A. Quilis,
\emph{Curve-flat functions and Lipschitz quotients},
arXiv:2603.20177.

\end{thebibliography}

\end{document}
